\section{Implementation}

The repository contains roughly 3{,}570 lines of our own code across three
languages, excluding the vendored \texttt{llama.cpp} checkout and generated
dependency files.

\subsection{Go controller}

The controller owns cluster state and the model process. It exposes
\code{POST /register}, \code{POST /heartbeat}, \code{GET /status},
\code{GET /pipeline} and \code{POST /infer}, and listens for discovery
broadcasts on UDP~50999. A worker is considered healthy if its last heartbeat
arrived within ten seconds.

Two implementation details are worth recording because both were bugs first.

The placement reader parses \texttt{llama.cpp}'s own load output. Running the
model server at log verbosity~4 makes it print one line per device of the form
\code{load\_tensors: <device> model buffer size = <N> MiB}, which the controller
captures into structured per-device figures. This is the only reason we noticed
the reversed tensor-split order described in Section~\ref{sec:background}: the
requested weights and the actual placement disagreed.

The process is started from the directory holding the \texttt{llama.cpp}
binaries. Renaming the project directory broke the build until we fixed this,
because the server resolves its shared libraries relative to its own location
and exited with status~127 before contacting any worker.

\subsection{Rust worker}

One agent binary serves Windows, Linux and macOS. It broadcasts presence,
registers, heartbeats every three seconds, and re-registers automatically if a
heartbeat fails, which is what allows the cluster to reassemble after the
controller restarts.

A single sampler runs every five seconds and its snapshot is shared by all three
message paths. This matters on Windows, where each sample would otherwise spawn
separate PowerShell and \code{nvidia-smi} processes for registration,
heartbeats and discovery.

\begin{table}[htbp]
\centering\small
\caption{Telemetry sources per platform. Every value these tools decline to
report is transmitted as \code{null}.}
\label{tab:sources}
\begin{tabular}{@{}llll@{}}
\toprule
 & \textbf{GPU} & \textbf{Memory} & \textbf{Power} \\
\midrule
Linux (NVIDIA)   & \code{nvidia-smi} & \code{/proc/meminfo} & \code{/sys/class/power\_supply} \\
Windows (NVIDIA) & \code{nvidia-smi} & CIM & CIM + WMI \\
macOS            & \code{system\_profiler} (name only) & \code{sysctl} + \code{vm\_stat} & \code{pmset} \\
\bottomrule
\end{tabular}
\end{table}

On macOS the chipset name is resolved once and cached, because
\code{system\_profiler} is slow and the value never changes, while memory and
power are sampled on every cycle.

\subsection{Dashboard}

The React dashboard polls \code{/status} and \code{/pipeline}. It shows the
topology, per-node measured hardware, the scheduler strategy with the exact
weights it produced, and the placement \texttt{llama.cpp} reported. Fields the
worker could not measure render as \textit{Unavailable}, never as zero.

\subsection{Tests}

Fifteen unit tests cover the parts where a silent error would be most damaging:
seven in Go and eight in Rust. They assert that unavailable values survive as
\code{null} through serialisation, that \code{[N/A]} fields are not coerced to
zero, that real \code{nvidia-smi} and \code{vm\_stat} output parses correctly,
that malformed output is rejected rather than half-read, that identical
measurements always yield identical tensor weights, that a missing measurement
produces the documented fallback instead of an invented weight, and that a
heartbeat carrying no telemetry does not erase the previous snapshot.

Continuous integration was deliberately left out of scope. All checks are run
locally: \code{go test}, \code{go vet}, \code{cargo test}, \code{cargo clippy}
with warnings as errors, and a TypeScript production build with zero lint
warnings.
